
\documentclass{article}

\input{setup}

\begin{document}

\CAPA{Trabalho Prático 1: Jogo da Vida}{BCC202 - Estruturas de Dados 1}{Fábio Tavares Pinto e Hugo Augusto Silva de Faria}{Pedro Henrique Lopes Silva}

\section{Introdução}

O Jogo da Vida, criado pelo matemático John Conway em 1970, é um exemplo de autômato celular que demonstra como sistemas simples podem gerar comportamentos complexos e imprevisíveis. Este jogo é baseado em uma grade bidimensional onde cada célula pode estar viva ou morta, e o estado de cada célula é atualizado em iterações discretas com base em regras específicas.

Neste trabalho, desenvolvemos um programa para simular o Jogo da Vida usando a linguagem C. O objetivo principal é criar e manipular um tabuleiro representado por uma matriz bidimensional de inteiros, onde cada célula pode assumir os valores 0 (morta) ou 1 (viva). O programa é projetado para realizar as seguintes operações:

\begin{itemize}
\item Alocação e Desalocação de Memória: Implementamos funções para alocar e liberar memória para o tabuleiro e sua matriz. Essas funções garantem que o espaço de memória seja gerido eficientemente durante a execução do programa.

\item Leitura e Impressão do Tabuleiro: Desenvolvemos funções para ler o estado inicial do tabuleiro a partir da entrada padrão e para imprimir o estado do tabuleiro na tela, permitindo a visualização das células vivas e mortas em cada iteração.

\item Evolução do Estado do Tabuleiro: O núcleo do programa é a implementação das regras do Jogo da Vida, que determinam como as células devem mudar de estado com base no número de vizinhos vivos. Implementamos uma função recursiva para evoluir o tabuleiro para a próxima geração e para executar um número específico de gerações. A recursividade simplifica a lógica do código, tornando-o mais intuitivo e alinhado com a natureza iterativa do Jogo da Vida.

\end{itemize}

Este relatório detalha a estrutura do programa, as funções implementadas e o funcionamento geral do sistema. Através da implementação deste simulador, buscamos entender melhor como sistemas complexos podem emergir a partir de regras simples e explorar a aplicação prática de autômatos celulares em computação. Além disso, o uso da recursividade na evolução do tabuleiro destaca a elegância e a eficiência de técnicas recursivas em problemas de transformação de estado.


\section{Arquivo de Cabeçalho (automato.h)}

O arquivo de cabeçalho automato.h define a interface para o módulo que implementa o Jogo da Vida. Ele contém as declarações de tipos de dados e funções que são utilizados para manipular o tabuleiro e executar as operações relacionadas ao jogo. A seguir, estão descritos os principais componentes e suas funções:

\subsection{Código}


\begin{lstlisting}[caption={},label={lst:cod1},language={[x86masm]Assembler}]
pragma once  // Garante que este arquivo de cabecalho sera incluido apenas uma vez no processo de compilacao

typedef struct {
    int n;            // A dimensao do tabuleiro (n x n)
    int **matriz;     // Ponteiro para uma matriz de inteiros que representa o tabuleiro
} Tabuleiro;

// Funcao para alocar memoria para o tabuleiro e a matriz
Tabuleiro* alocarReticulado(int n);

// Funcao para desalocar a memoria do tabuleiro
void desalocarReticulado(Tabuleiro **tab);

// Funcao para ler o tabuleiro
void LeituraReticulado(Tabuleiro *tab);

// Funcao para evoluir o estado do tabuleiro para a geracao N de acordo com as regras do jogo da vida
Tabuleiro* evoluirReticulado(Tabuleiro *tab, int n,int geracoes);

// Funcao para imprimir o estado do tabuleiro na tela
void imprimeReticulado(Tabuleiro *tab);
 \end{lstlisting}

 \subsection{Estrutura 'Tabuleiro'}

A estrutura Tabuleiro é definida para representar o estado do tabuleiro do Jogo da Vida. Ela contém:

\begin{itemize}
    \item 'int n': Um inteiro que representa a dimensão do tabuleiro, sendo 'n x n'.
    \item 'int **matriz': Um ponteiro para um ponteiro de inteiros, que representa a matriz bidimensional do tabuleiro. Cada célula da matriz pode ser 0 (morta) ou 1 (viva).
\end{itemize}

 \subsection{Funções Declaradas}

 \begin{itemize}

 \item Tabuleiro* alocarReticulado(int n): Esta função aloca memória para um novo tabuleiro com a dimensão n x n. Ela retorna um ponteiro para a estrutura Tabuleiro recém-alocada. Se a alocação falhar, retorna NULL.

 \item void desalocarReticulado(Tabuleiro **tab):Esta função libera a memória alocada para um tabuleiro. Ela percorre cada linha da matriz, liberando a memória associada, depois libera a matriz e, finalmente, libera a estrutura Tabuleiro, garantindo que todos os recursos de memória sejam corretamente desalocados.

\item void LeituraReticulado(Tabuleiro *tab): Lê o estado inicial do tabuleiro a partir da entrada padrão. Preenche a matriz do tabuleiro com os valores fornecidos, onde cada célula pode ser 0 ou 1.

\item Tabuleiro* evoluirReticulado(Tabuleiro *tab, int n,int geracoes): Esta função recursiva atualiza o estado do tabuleiro para a próxima geração com base nas regras do Jogo da Vida de Conway. Ela aloca um novo tabuleiro, aplica as regras de vizinhança para cada célula, desaloca o tabuleiro antigo e chama a si mesma para as gerações restantes, retornando o tabuleiro final após todas as gerações.

\item void imprimeReticulado(Tabuleiro *tab): Imprime o estado atual do tabuleiro na tela. Cada célula é exibida com um espaço entre os valores, e cada linha do tabuleiro é impressa em uma nova linha.

 \end{itemize}


\section{Arquivo de Implementação (automato.c)}

O arquivo de implementação automato.c contém a definição das funções declaradas no arquivo de cabeçalho automato.h. Essas funções são responsáveis pela alocação, desalocação, manipulação e evolução do tabuleiro do Jogo da Vida. A seguir, estão descritas as principais funções e sua implementação:

 \subsection{Função alocarReticulado}
\begin{lstlisting}[caption={},label={lst:cod1},language={[x86masm]Assembler}]
Tabuleiro *alocarReticulado(int n)
{
    Tabuleiro *tab = (Tabuleiro *)malloc(sizeof(Tabuleiro));
    if (tab == NULL)
    {
        printf("Erro de memoria\n");
        return NULL;
    }

    tab->n = n;

    tab->matriz = (int **)malloc(n * sizeof(int *));
    if (tab->matriz == NULL)
    {
        printf("Erro de memoria\n");
        free(tab);
        return NULL;
    }

    for (int i = 0; i < n; i++)
    {
        tab->matriz[i] = (int *)malloc(n * sizeof(int));
        if (tab->matriz[i] == NULL)
        {
            for (int j = 0; j < i; j++)
            {
                free(tab->matriz[j]);
            }
            free(tab->matriz);
            free(tab);
            printf("Erro de memoria\n");
            return NULL;
        }
    }

    return tab;
}
 \end{lstlisting}

 \par A função alocarReticulado é responsável por alocar memória para o struct Tabuleiro e inicializar sua matriz. Inicialmente, a função aloca memória para o struct Tabuleiro e verifica se a alocação foi bem-sucedida. Se houver um erro na alocação, a função imprime uma mensagem de erro e retorna NULL. Em seguida, a dimensão do tabuleiro (n) é inicializada. A função então aloca memória para a matriz, que é uma matriz 2D de inteiros. Para cada linha da matriz, a função tenta alocar memória e, se houver falha em qualquer ponto, libera toda a memória já alocada e retorna NULL, imprimindo uma mensagem de erro. Finalmente, a função retorna um ponteiro para o struct Tabuleiro alocado.

\subsection{Função desalocarReticulado}

\begin{lstlisting}[caption={},label={lst:cod1},language={[x86masm]Assembler}]
void desalocarReticulado(Tabuleiro **tab)
{
    if (tab == NULL || *tab == NULL || (*tab)->matriz == NULL)
    {
        return;
    }

    for (int i = 0; i < (*tab)->n; i++)
    {
        free((*tab)->matriz[i]);
    }

    free((*tab)->matriz);
    free(*tab);
    *tab = NULL;
}

 \end{lstlisting}

 \par A função desalocarReticulado cuida da liberação da memória alocada para o struct Tabuleiro e sua matriz. A função começa verificando se o ponteiro do tabuleiro ou a matriz são nulos. Se qualquer um deles for nulo, a função retorna imediatamente. Caso contrário, a função libera a memória alocada para cada linha da matriz. Após liberar todas as linhas, a função libera a memória da própria matriz e do struct Tabuleiro. Para evitar acessos posteriores a um ponteiro inválido, a função define o ponteiro do tabuleiro como NULL. Isso garante que qualquer tentativa de acesso ao tabuleiro após sua desalocação resulte em um ponteiro nulo, ajudando a prevenir erros de memória.

 \subsection{Função LeituraReticulado}

\begin{lstlisting}[caption={},label={lst:cod1},language={[x86masm]Assembler}]
void LeituraReticulado(Tabuleiro *tab)
{
    for (int i = 0; i < tab->n; i++)
    {
        for (int j = 0; j < tab->n; j++)
        {
            scanf("%d", &tab->matriz[i][j]);
        }
    }
}


 \end{lstlisting}
 
\par A função LeituraReticulado é responsável por ler o estado inicial do tabuleiro a partir da entrada padrão. Ela percorre todas as células da matriz do tabuleiro e usa a função scanf para ler o valor de cada célula. Essa função é essencial para inicializar o tabuleiro com os valores que representam o estado inicial de cada célula (viva ou morta).


\subsection{Função evoluirReticulado}
\begin{lstlisting}[caption={},label={lst:cod1},language={[x86masm]Assembler}]
Tabuleiro *evoluirReticulado(Tabuleiro *tab, int n, int geracoes)
{
    // Caso base: se nao ha mais geracoes para evoluir, retorne o tabuleiro atual.
    if (geracoes == 0)
    {
        return tab; // Se o numero de geracoes restantes e 0, retornamos o tabuleiro atual, que representa o estado apos todas as geracoes solicitadas.
    }

    // Evolui o tabuleiro para a proxima geracao.
    // Cria um novo tabuleiro para armazenar o estado da proxima geracao.
    Tabuleiro *novoTabuleiro = alocarReticulado(n);

    // Verifica se a alocacao foi bem-sucedida.
    if (novoTabuleiro == NULL)
    {
        return NULL; // Se a alocacao falhar, retorna NULL. Isso indica um erro na alocacao de memoria.
    }

    // Aplica as regras do Jogo da Vida para definir o proximo estado da celula (i, j).
    // Percorre cada celula no tabuleiro atual.
    for (int i = 0; i < n; i++)
    {
        for (int j = 0; j < n; j++)
        {
            int vizinhosVivos = 0; // Conta o numero de vizinhos vivos da celula (i, j).

            // Verifica todas as celulas ao redor da celula atual (i, j), incluindo diagonais.
            for (int di = -1; di <= 1; di++)
            {
                for (int dj = -1; dj <= 1; dj++)
                {
                    // Ignora a propria celula (i, j).
                    if (di == 0 && dj == 0)
                        continue;

                    int ni = i + di; // Calcula a linha do vizinho.
                    int nj = j + dj; // Calcula a coluna do vizinho.

                    // Verifica se o vizinho esta dentro dos limites do tabuleiro.
                    if (ni >= 0 && ni < n && nj >= 0 && nj < n)
                    {
                        // Adiciona o valor da celula vizinha ao contador de vizinhos vivos.
                        vizinhosVivos += tab->matriz[ni][nj];
                    }
                }
            }

            // Aplica as regras do Jogo da Vida para atualizar o estado da celula (i, j).
            if (tab->matriz[i][j] == 1) // Se a celula esta viva.
            {
                // A celula continua viva se tiver 2 ou 3 vizinhos vivos; caso contrario, morre.
                if (vizinhosVivos == 2 || vizinhosVivos == 3)
                    novoTabuleiro->matriz[i][j] = 1; // Continua viva.
                else
                    novoTabuleiro->matriz[i][j] = 0; // Morre.
            }
            else // Se a celula esta morta.
            {
                // A celula se torna viva se tiver exatamente 3 vizinhos vivos; caso contrario, permanece morta.
                if (vizinhosVivos == 3)
                    novoTabuleiro->matriz[i][j] = 1; // Torna-se viva.
                else
                    novoTabuleiro->matriz[i][j] = 0; // Permanece morta.
            }
        }
    }

    // Desaloca o tabuleiro antigo para liberar memoria.
    desalocarReticulado(&tab);

    // Chama a funcao recursivamente para o proximo numero de geracoes.
    // Atualiza o tabuleiro para a nova geracao e decrementa o numero de geracoes restantes.
    return evoluirReticulado(novoTabuleiro, n, geracoes - 1);
}
 \end{lstlisting}

\par A função evoluirReticulado é responsável por evoluir um tabuleiro do Jogo da Vida de Conway através de múltiplas gerações. O propósito desta função é calcular o estado do tabuleiro após um número específico de gerações, utilizando recursão para alcançar o resultado desejado.

\par A função recebe três parâmetros: um ponteiro para um Tabuleiro (tab), o tamanho do tabuleiro (n), e o número de gerações a serem aplicadas (geracoes). O processo inicia verificando se o número de gerações (geracoes) é igual a zero. Se for, a função retorna o tabuleiro atual sem alterações, pois não há mais gerações a serem aplicadas.

\par Se ainda houver gerações a serem processadas, a função continua a execução criando um novo tabuleiro (novoTabuleiro) que armazenará o estado da próxima geração. A função alocarReticulado é utilizada para alocar memória para este novo tabuleiro. Caso a alocação falhe e novoTabuleiro seja NULL, a função retorna NULL para indicar um erro de alocação de memória.

\par Com o novo tabuleiro alocado com sucesso, a função procede para aplicar as regras do Jogo da Vida de Conway. Para isso, a função itera sobre cada célula do tabuleiro atual. Para cada célula na posição (i, j), a função conta o número de vizinhos vivos. Esta contagem é feita percorrendo todas as células ao redor da célula atual, incluindo as diagonais, e somando o valor das células vizinhas.

\par Após calcular o número de vizinhos vivos, a função aplica as regras do Jogo da Vida para determinar o próximo estado da célula (i, j) no novo tabuleiro. Se a célula está viva e possui 2 ou 3 vizinhos vivos, ela permanece viva; caso contrário, ela morre. Se a célula está morta e possui exatamente 3 vizinhos vivos, ela se torna viva; caso contrário, permanece morta.

\par Com o novo tabuleiro atualizado, a função desaloca o tabuleiro antigo utilizando a função desalocarReticulado para liberar a memória previamente alocada. Em seguida, a função chama a si mesma recursivamente para processar a próxima geração, passando o novoTabuleiro como o tabuleiro atual e decrementando o número de gerações restantes.

\par Finalmente, a função retorna o resultado da chamada recursiva, que é o tabuleiro após o número total de gerações especificado. Este retorno representa o estado final do tabuleiro após todas as gerações.

\subsection{Função imprimeReticulado}

\begin{lstlisting}[caption={},label={lst:cod1},language={[x86masm]Assembler}]
void imprimeReticulado(Tabuleiro *tab)
{
    printf("\n");

    for (int i = 0; i < tab->n; i++)
    {
        for (int j = 0; j < tab->n; j++)
        {
            printf("%d ", tab->matriz[i][j]);
        }
        printf("\n");
    }
}


\end{lstlisting}

\par A função imprimeReticulado imprime o estado atual do tabuleiro na tela. Ela percorre todas as células do tabuleiro e usa a função printf para imprimir o valor de cada célula, seguido por um espaço. Após imprimir todas as células de uma linha, a função imprime uma nova linha. Isso é repetido para todas as linhas do tabuleiro, resultando em uma representação visual do estado atual do tabuleiro no terminal.



\section{Arquivo de Implementação (tp.c)}

O arquivo de implementação tp.c contém o código principal que utiliza as funções declaradas no arquivo de cabeçalho automato.h e implementadas em automato.c. Este arquivo orquestra a execução do programa, gerencia a entrada e saída de dados, e controla a evolução do tabuleiro ao longo das gerações. A seguir, está descrito o funcionamento das principais seções e funções deste arquivo:

 \subsection{Inclusão de Cabeçalhos}


\begin{lstlisting}[caption={},label={lst:cod1},language={[x86masm]Assembler}]
include <stdio.h>
include <stdlib.h>
include "automato.h"

 \end{lstlisting}

\par Inclui as bibliotecas necessárias para a entrada e saída padrão e alocação de memória, além do arquivo de cabeçalho automato.h que contém as declarações das funções e definições do Tabuleiro.

\subsection{Função main}

\begin{lstlisting}[caption={},label={lst:cod1},language={[x86masm]Assembler}]
int main()
{
    int n, geracao;

    scanf("%d %d", &n, &geracoes);

 \end{lstlisting}

 \par Declaração das variáveis n (dimensão do tabuleiro) e geracao (número de gerações a serem evoluídas). Utiliza scanf para ler a dimensão do tabuleiro e o número de gerações a partir da entrada padrão.

\subsection{Chamadas das Funções}

\begin{lstlisting}[caption={},label={lst:cod1},language={[x86masm]Assembler}]
Tabuleiro *tab = alocarReticulado(n);

LeituraReticulado(tab);

tab = evoluirReticulado(tab, n, geracoes);

imprimeReticulado(tab);

desalocarReticulado(&tab);

\end{lstlisting}

\begin{itemize}

\item  Chama a função alocarReticulado para alocar memória para o tabuleiro com a dimensão especificada.

\item Chama a função LeituraReticulado para ler o estado inicial do tabuleiro a partir da entrada padrão.

\item Chama a função evoluirReticulado para evoluir o tabuleiro pelo número de gerações especificado. O tabuleiro evoluído substitui o original.

\item Chama a função imprimeReticulado para imprimir o estado final do tabuleiro.

\item Chama a função desalocarReticulado para liberar a memória alocada para o tabuleiro.

\end{itemize}

\subsection{Fim da Função main}

\begin{lstlisting}[caption={},label={lst:cod1},language={[x86masm]Assembler}]
return 0;
}

 \end{lstlisting}

 \par Retorna 0, indicando que o programa terminou com sucesso.

\section{Impressões gerais}

\par Durante a implementação deste trabalho, gostei particularmente do desafio de gerenciar a alocação e desalocação de memória, o que me permitiu aprofundar meu conhecimento sobre a manipulação de matrizes dinâmicas em C.
\par Foi motivador ver o tabuleiro evoluir conforme as regras do Jogo da Vida, destacando a beleza de como regras simples podem resultar em comportamentos complexos. No entanto, encontrei dificuldades na depuração de erros de memória, o que foi inicialmente frustrante, mas extremamente educativo.
\par Além disso, implementar a recursividade na função de evolução do tabuleiro foi um desafio interessante. A abordagem recursiva não só simplificou a lógica do código, mas também aprimorou meu entendimento sobre como a recursividade pode ser aplicada em algoritmos de transformação de estado.
\par Este projeto me proporcionou um entendimento mais sólido de autômatos celulares e melhorou minhas habilidades em C, especialmente no que diz respeito ao gerenciamento eficiente de recursos e à aplicação de técnicas recursivas.

\section{Análise}

\par Vamos analisar o tempo de execução do código e se teve algum vazamento de memória.

\subsection{Tempo de Execução}
\par Testamos o tempo de execução da entrada do teste 20.in. O resultado obtido foi:

\begin{lstlisting}[caption={},label={lst:cod1},language={[x86masm]Assembler}]
real    0m4.859s
user    0m0.003s
sys     0m0.001s

 \end{lstlisting}

\subsection{Valgrind}
\par Utilizamos o Valgrind para verificar se teve algum vazamento de memória, também com a entrada do teste 20.in. O resultado obtido foi:

\begin{lstlisting}[caption={},label={lst:cod1},language={[x86masm]Assembler}]
==159134== HEAP SUMMARY:
==159134==     in use at exit: 0 bytes in 0 blocks
==159134==   total heap usage: 46 allocs, 46 frees, 5,600 bytes allocated
==159134== 
==159134== All heap blocks were freed -- no leaks are possible
==159134== 
==159134== ERROR SUMMARY: 0 errors from 0 contexts (suppressed: 0 from 0)

 \end{lstlisting}





\section{Conclusão}
\par Após a implementação e execução do nosso programa, podemos concluir que o desempenho e a eficiência do código foram satisfatórios. Utilizamos o comando time para medir o tempo de execução do nosso programa e observamos que o tempo de execução foi excelente, evidenciando a eficiência das funções implementadas.

\par Além disso, utilizamos a ferramenta valgrind para verificar possíveis vazamentos de memória durante a execução do programa. Os resultados mostraram que não houve vazamentos de memória, indicando que a alocação e desalocação de memória foram realizadas corretamente em todas as funções do código.

\par O programa cumpriu com sucesso todos os requisitos propostos, evoluindo corretamente o estado do tabuleiro conforme as regras do Jogo da Vida de Conway. A impressão do tabuleiro nas diversas gerações também foi realizada conforme esperado, proporcionando uma visualização clara da evolução das células ao longo do tempo. A implementação da recursividade na função de evolução do tabuleiro também se mostrou eficaz, simplificando a lógica do código e demonstrando a aplicabilidade de técnicas recursivas em algoritmos de transformação de estado.

\par Em resumo, o projeto alcançou seus objetivos com um código eficiente e sem problemas de memória, demonstrando um bom entendimento e aplicação das técnicas de alocação dinâmica de memória, manipulação de matrizes e recursividade em C. O resultado final foi obtido com sucesso, validando a eficácia das soluções implementadas.


\end{document}